\documentclass[9pt,a4paper]{article}
\usepackage[spanish,es-nodecimaldot]{babel}
\usepackage[utf8]{inputenc}
\usepackage[T1]{fontenc}
\usepackage{lmodern}
\usepackage[landscape,margin=0.5cm]{geometry}
\usepackage{multicol}
\usepackage{amsmath,amssymb}
\usepackage[table]{xcolor}
\usepackage{booktabs,array,multirow}
\usepackage{enumitem}
\usepackage[normalem]{ulem}

\pagestyle{empty}
\setlength{\parindent}{0pt}
\setlength{\parskip}{1pt}
\setlength{\columnsep}{8pt}
\setlength{\columnseprule}{0.3pt}

\AtBeginDocument{%
  \setlength{\abovedisplayskip}{1pt}%
  \setlength{\belowdisplayskip}{1pt}%
  \setlength{\abovedisplayshortskip}{0.5pt}%
  \setlength{\belowdisplayshortskip}{0.5pt}%
}

\newcommand{\bodysize}{\fontsize{7.5}{8.7}\selectfont}
\newcommand{\notesize}{\fontsize{6.4}{7.4}\selectfont}
\newcommand{\subsize}{\fontsize{8.6}{10}\selectfont\bfseries}
\newcommand{\unitsize}{\fontsize{9.8}{11.3}\selectfont\bfseries}

\newcommand{\unithead}[2]{%
  \par\vspace{2pt}
  \noindent\colorbox{#1!22}{\textcolor{#1!65!black}{\unitsize\ #2}}%
  \par\vspace{1pt}\hrule height 0.6pt \color{#1!55!black}\vspace{1.5pt}\color{black}%
}
\newcommand{\topic}[2]{\subsize\textcolor{#1!60!black}{#2}\par\vspace{0.5pt}\bodysize}
\newcommand{\ejemplo}[1]{{\notesize\textbf{Ej:} #1\par}\bodysize}

\bodysize

\begin{document}
\bodysize

\begin{multicols}{5}

%======================= UNIDAD 4 =======================
\unithead{violet}{UNIDAD 4 --- No Param.}

\topic{violet}{Prueba del Signo (1 muestra o apareadas)}
Alternativa no paramétrica a t-Unimuestral. Reemplazar cada valor por $+$/$-$ según sea mayor/menor a la media o mediana (se ignora si es igual); prob. $p=0.5$ siempre. $x$=cant. de signos $+$, $n$=total de signos. Se calcula $P(X\geq x)$ (dbinom) y si es $\leq\alpha$ se rechaza $H_0$.

\ejemplo{Octanaje de nafta, $n=15$ ($\mu_0=98$): 1 empate se descarta $\to n=14$, $x=12$ signos $+$. $H_0:\tilde\mu=98$, $H_1:\tilde\mu>98$. $p$-valor$=P(X\geq12)=0.0065<\alpha=0.01$ $\Rightarrow$ \textbf{Rechaza $H_0$}.}

\ejemplo{Accidentes antes/después (pareada), $n=10$ plantas: $x=9$ signos $+$ (antes$>$después). $H_0:\tilde\mu_D=0$, $H_1:\tilde\mu_D>0$. $p$-valor$=P(X\geq9)=0.0107<\alpha=0.05$ $\Rightarrow$ \textbf{Rechaza $H_0$} (programa eficaz).}

\topic{violet}{Prueba U (Wilcoxon, 2 muestras indep.)}
Alternativa no paramétrica a t-bimuestral. Unir y ordenar ambas muestras; asignar rango (posición, promedio si hay empate); sumar $R_1,R_2$ y elegir el $R_i$ menor.
\[
U_i=n_1n_2+\frac{n_i(n_i+1)}{2}-R_i
\]
Si $n_1,n_2>8$: $U_i$ aprox. normal (2 colas).
\[
z=\frac{U_i-\mu_{U_i}}{\sigma_{U_i}}
\]
\[
\mu_{U1}=\frac{n_1n_2}{2}
\]
\[
\sigma_{U1}=\sqrt{\frac{n_1n_2(n_1+n_2+1)}{12}}
\]

\topic{violet}{Prueba H (Kruskal-Wallis, $k\geq3$ indep.)}
Alternativa no paramétrica a ANOVA. $H_0:\mu_1=\cdots=\mu_k$. Unir las $k$ muestras, ordenar creciente, rango (promedio en empates), sumar $R_i$ por grupo.
\[
H=\frac{12}{n(n+1)}\sum_{i=1}^{k}\frac{R_i^2}{n_i}-3(n+1)
\]
Válido si $n_i>5$. $H\sim\chi^2_{k-1}$. $n$=suma de todos los tamaños muestrales.

\ejemplo{Corrosión en cables, 3 métodos ($n_1{=}6,n_2{=}7,n_3{=}5$, $N=18$):
\begin{center}
{\notesize
\setlength{\tabcolsep}{3pt}
\begin{tabular}{@{}l c c@{}}
\toprule
Método & $n_i$ & $R_i$ \\
\midrule
A & 6 & 84 \\
B & 7 & 55.5 \\
C & 5 & 31.5 \\
\bottomrule
\end{tabular}
}
\end{center}
$H=6.66>\chi^2_{crit}(2)=5.991$ $\Rightarrow$ \textbf{Rechaza $H_0$} (A el menos eficaz, C el mejor).}

\topic{violet}{Pruebas de aleatoriedad (corridas)}
$H_0$: secuencia aleatoria. $H_1$: no aleatoria (bil.) / demasiadas corridas (unil. der.) / muy pocas corridas (unil. izq.). Corrida=letras iguales consecutivas. Binario: cadena a/b directa. Numérico: $a$ si $>$mediana, $b$ si $<$, se omiten los iguales, se mantiene el orden.
\[
\mu_u=\frac{2n_1n_2}{n_1+n_2}+1
\]
\[
\sigma_u=\sqrt{\frac{2n_1n_2(2n_1n_2-n_1-n_2)}{(n_1+n_2)^2(n_1+n_2-1)}}
\]
Si $n_1,n_2>10$: $Z_{calc}=(u-\mu_u)/\sigma_u$. $u$ grande=alternancia excesiva; $u$ chico=valores agrupados; $u\approx\mu_u$=compatible con aleatoriedad.

%======================= UNIDAD 5 =======================
\unithead{orange}{UNIDAD 5 --- Regresión}

\topic{orange}{Método de los mínimos cuadrados}
$\hat{y}=a+bx$ estima $y=\alpha+\beta x$; $e_i=y_i-\hat y_i$ (dif. verticales). Minimizar $\sum[y_i-(a+bx_i)]^2$.
{\notesize
\[
\sum y_i=an+b\sum x_i
\]
\[
\sum y_ix_i=a\sum x_i+b\sum x_i^2
\]
\[
\begin{bmatrix} n & \sum x_i \\ \sum x_i & \sum x_i^2 \end{bmatrix}
\begin{pmatrix} a\\b \end{pmatrix} =
\begin{bmatrix} \sum y_i \\ \sum y_ix_i \end{bmatrix}
\]
}
$a$=\texttt{intercept(x,y)}, $b$=\texttt{slope(x,y)}.

\topic{orange}{Correlación}
\[
r=\frac{s_{xy}}{s_x\cdot s_y}=\frac{\sum(x_i-\bar x)(y_i-\bar y)}{\sqrt{\sum(x_i-\bar x)^2\cdot\sum(y_i-\bar y)^2}}
\]
$-1\leq r\leq1$. $r=1$: recta pend. positiva perfecta. $r=-1$: recta pend. negativa perfecta. $r=0$: sin asociación lineal (puede haber relación no lineal). $r$ mide relación entre 2 variables; \textbf{no} permite explicar una en función de la otra (no es causalidad).

\topic{orange}{Inferencias basadas en mín. cuadrados}
\[
b=\frac{S_{xy}}{S_{xx}} \qquad a=\bar Y-b\bar X
\]
{\notesize
\[
s_e^2=\frac{S_{xx}S_{yy}-S_{xy}^2}{n(n-2)S_{xx}}\ (\nu=n-2) \qquad s_e=\sqrt{s_e^2}
\]
}
\textbf{IC:}
\[
a\pm t_{\alpha/2}s_e\sqrt{\frac{nS_{xx}}{S_{xx}+(n\bar x)^2}}
\]
\[
b\pm t_{\alpha/2}s_e\sqrt{\frac{n}{S_{xx}}}
\]
$H_0:\beta=0$ vs $H_1:\beta\neq0$ (¿$X$ influye linealmente en $Y$? si $\beta=0$, $Y$ no depende de $X$). Estimación puntual: $\hat y_0=a+bx_0$.

\topic{orange}{Regresión curvilínea (linealizar)}
\textbf{Exponencial} $y=\alpha\beta^x\Rightarrow y^*=\log_{10}y=a+bx$, $\alpha=10^a,\beta=10^b$.
\textbf{Potencial} $y=\alpha x^\beta\Rightarrow y^*=\log_{10}y,\ x^*=\log_{10}x: y^*=a+Bx^*$.
\textbf{Recíproca} $y=\dfrac{1}{\alpha+\beta x}\Rightarrow y^*=\dfrac{1}{y}=\alpha+\beta x$.

\topic{orange}{Ajuste polinómico}
Grado $p$, coef. $b_0,\ldots,b_p$; minimizar $\sum[y_i-(b_0+b_1x_i+\cdots+b_px_i^p)]^2$ y armar sistema de $p+1$ ecuaciones normales (matriz de sumas de potencias de $x$).

\topic{orange}{Método de la varianza residual}
\[
\sigma_{res}^2=\frac{\sum(y_i-\hat y_i)^2}{n-(p+1)}
\]
Se prueba grado 1 (recta), 2 (parábola), 3,\ldots y se elige el grado donde se observa la mayor caída de $\sigma_{res}^2$; subir más el grado aumenta el riesgo de sobreajuste.

\ejemplo{$n=10$ mediciones:
\begin{center}
{\notesize
\setlength{\tabcolsep}{3pt}
\begin{tabular}{@{}c c c c@{}}
\toprule
$p$ & SCE & gl & $s^2$ \\
\midrule
1 (Lineal) & 57.655 & 8 & 7.207 \\
2 (Cuadr.) & 53.059 & 7 & 7.580 \\
3 (Cúbico) & 6.124 & 6 & \textbf{1.021} \\
4 (Cuárt.) & 4.959 & 5 & 0.992 \\
\bottomrule
\end{tabular}
}
\end{center}
Mayor caída en $p=3$ (salto de 7.58 a 1.02) $\Rightarrow$ se elige polinomio cúbico.}

\topic{orange}{Regresión múltiple}
Incorpora $x_1,\ldots,x_k$ variables independientes para explicar $Y$; se estima un plano (o hiperplano) de regresión en vez de una recta.

%======================= UNIDAD 6 =======================
\unithead{green}{UNIDAD 6 --- ANOVA}

\topic{green}{Diseño completamente aleatorio (1 vía)}
$H_0:\mu_1=\cdots=\mu_k$; $H_1$: al menos una distinta. $k$ tratamientos, $n$ obs. c/u. $\alpha_i=\bar y_{i.}-\bar y_{..}$ (si $H_0$ cierta, todos $\alpha_i=0$).
\[
C=\frac{T_{..}^2}{kn} \qquad SST=\sum\sum y_{ij}^2-C
\]
{\notesize
\[
SS(Tr)=\frac{1}{n}\sum T_{i.}^2-C=n\sum(\bar y_{i.}-\bar y_{..})^2
\]
}
\[
SSE=SST-SS(Tr)
\]
\[
F=\frac{SS(Tr)/(k-1)}{SSE/[k(n-1)]}
\]
\begin{center}
{\notesize
\begin{tabular}{@{}l c c c@{}}
\toprule
Fuente & gl & SC & MC \\
\midrule
Tratam. & $k-1$ & $SS(Tr)$ & $MS(Tr)$ \\
Error & $k(n-1)$ & $SSE$ & $MSE$ \\
Total & $nk-1$ & $SST$ & \\
\bottomrule
\end{tabular}
}
\end{center}

\ejemplo{4 laboratorios de estaño, $k=4,n=12,N=48$:
\begin{center}
{\notesize
\setlength{\tabcolsep}{3pt}
\begin{tabular}{@{}l c c c c@{}}
\toprule
Fuente & gl & SC & MC & $F$ \\
\midrule
Labs. & 3 & 0.0130 & 0.00433 & \textbf{2.87} \\
Error & 44 & 0.0679 & 0.00154 & \\
\midrule
Total & 47 & 0.0809 & & \\
\bottomrule
\end{tabular}
}
\end{center}
$F_{crit}(3,44)=2.82$. Como $2.87>2.82$ $\Rightarrow$ \textbf{Rechaza $H_0$} (evidencia marginal, sugiere aumentar $n$).}

\topic{green}{Diseño en bloques aleatorios}
$y_{ij}=\mu+\alpha_i+\beta_j+\epsilon_{ij}$, $\sum\alpha_i=0,\sum\beta_j=0$; $\epsilon_{ij}\sim iid\,N(0,\sigma^2)$. $a$=nº tratamientos, $b$=nº bloques.
\[
C=\frac{T_{..}^2}{ab} \qquad SST=\sum\sum y_{ij}^2-C
\]
\[
SS(Tr)=\frac{1}{b}\sum T_{i.}^2-C\ (a{-}1\text{ gl})
\]
\[
SS(Bl)=\frac{1}{a}\sum T_{.j}^2-C\ (b{-}1\text{ gl})
\]
{\notesize
\[
SSE=SST-SS(Tr)-SS(Bl)
\]
\[
[(a{-}1)(b{-}1)\text{ gl}]
\]
}
\[
F_{Tr}=\frac{MS(Tr)}{MSE} \qquad F_{Bl}=\frac{MS(Bl)}{MSE}
\]
\begin{center}
{\notesize
\begin{tabular}{@{}l c c@{}}
\toprule
Fuente & gl & SC \\
\midrule
Tratam. & $a-1$ & $SS(Tr)$ \\
Bloques & $b-1$ & $SS(Bl)$ \\
Error & $(a{-}1)(b{-}1)$ & $SSE$ \\
Total & $ab-1$ & $SST$ \\
\bottomrule
\end{tabular}
}
\end{center}

\ejemplo{Detergentes ($a=3$) $\times$ Lavarropas ($b=3$):
\begin{center}
{\notesize
\setlength{\tabcolsep}{3pt}
\begin{tabular}{@{}l c c c@{}}
\toprule
Fuente & gl & SC & $F_{calc}$ \\
\midrule
Deterg. & 2 & 600 & \textbf{120} \\
Lavarr. & 2 & 150 & \textbf{30} \\
Error & 4 & 10 & \\
\midrule
Total & 8 & 760 & \\
\bottomrule
\end{tabular}
}
\end{center}
$F_{crit}(2,4;0.01)=18$. Como $120>18$ y $30>18$ $\Rightarrow$ \textbf{Rechaza ambas $H_0$} (detergentes y lavarropas influyen).}

\topic{green}{Tamaños muestrales distintos ($n_i$)}
Cambios respecto al diseño 1 vía: usar $n_i$ (no $n$) en $SST,SS(Tr)$ y $T$; usar $N=\sum n_i$ (no $k\cdot n$) en $C$.
\[
C=\frac{T_{..}^2}{N} \qquad SST=\sum\sum y_{ij}^2-C
\]
{\notesize
\[
SS(Tr)=\sum\frac{T_{i.}^2}{n_i}-C
\]
\[
SSE=SST-SS(Tr)
\]
}
gl: Tratamientos $k-1$, Error $N-k$, Total $N-1$.

\ejemplo{Aflatoxina en 2 marcas, $k=2,n_1=8,n_2=6,N=14$:
\begin{center}
{\notesize
\setlength{\tabcolsep}{3pt}
\begin{tabular}{@{}l c c c@{}}
\toprule
Fuente & gl & SC & $F_{calc}$ \\
\midrule
Marcas & 1 & 11.734 & \textbf{1.047} \\
Error & 12 & 134.515 & \\
\midrule
Total & 13 & 146.249 & \\
\bottomrule
\end{tabular}
}
\end{center}
$F_{crit}(1,12)=4.747$. Como $1.047<4.747$ $\Rightarrow$ \textbf{No se rechaza $H_0$} (las marcas no difieren).}

\end{multicols}
\end{document}
